\chapter{Etude conceptuelle}

\section*{Modélisation}
\addcontentsline{toc}{section}{Modélisation}
L'étude conceptuelle frontend s'est appuyée sur la définition des parcours utilisateurs et des flux de navigation.

Les éléments principaux modélisés sont :
\begin{itemize}
	\item architecture des routes (landing, login, dashboard, pages spécialisées) ;
	\item segmentation des fonctionnalités par rôle ;
	\item composants génériques (ResourcePage, notifications, profils) ;
	\item scénarios de navigation sécurisée via \textit{ProtectedRoute}.
\end{itemize}

\section*{Répartition des tâches conceptuelles (Lot B)}
\addcontentsline{toc}{section}{Répartition des tâches conceptuelles (Lot B)}
\begin{tabularx}{\textwidth}{lX}
\toprule
\textbf{Axe} & \textbf{Tâches Lot B} \\
\midrule
UX/UI & Conception des écrans, homogénéisation de la charte visuelle et organisation des composants. \\
Parcours métier & Scénarios enseignant/étudiant : cours, devoirs, notes, absences, demandes, notifications. \\
Intégration API & Mapping des champs formulaires avec les endpoints backend et gestion des réponses. \\
\bottomrule
\end{tabularx}
